% Integration fragment, 2026-10-08.  Does not modify the locked second seed.
\subsection{The second seed on stable specified-prefix families}
\label{sec:second-seed-stable-prefix}

The numerical seed used here is exactly the proved $J_*=9/500$ of
Proposition~\ref{prop:replacement-second-seed}, with
\[
 s_*=11/5,\quad s_0=227/50,\quad
 (\phi_1,\phi_2,\phi_3)=(353/100,29/10,61/25).
\]
Its proof is now extended to the stable factorization branches used
in Section~\ref{sec:generalized-one-step-double-sieve}.  No old
correction table is used and no new numerical integration is needed.
The extension keeps the repaired $\Gamma_5$ kernel, all other
ordinary kernels, and the separate unit masses in $\Gamma_{20},
\Gamma_{21}$ without change.

Write the specified squarefree prefix as $B=p_1\cdots p_r$, with
decreasing exponents $t_1>\cdots>t_r>0$, and put
\[
 T=\sum_i t_i,\qquad Q=N^\nu,\quad R=Q/B=N^{L+o(1)},\quad
 L=\nu-T>0,\qquad \alpha=(1-T)/L\ge2.
\]
All prefix coefficients are finite nonnegative constant combinations
of marked prime-box indicators.  The prime exponents and $L$ have
fixed positive lower bounds on every box.  As usual all box,
factorization, and level inequalities have strict margins before
taking their endpoint limits.

\paragraph{Precisely which factorization data are retained.}
In the enhanced case $1/2<\nu<153/256$, require the full specified
square conditions
\[
 t_1+\cdots+t_{j-1}+2t_j<\nu\qquad(1\le j\le r).
\]
Choose one of the already proved cap branches, below its cap and
above every gate: the largest-prime branch; a selected-subproduct
branch using specified slots and their remaining-square gates; or
the repaired $W$ branch with its first three slots or a retained
bound on its third slot.  The chosen branch must persist on adding
primes of exponent below a specified ceiling $\chi\le t_r$.
For branches determined by existing slots one can take $\chi=t_r$.
If there are fewer than three specified slots and the $W$ branch
uses only the assertion that a future third exponent is at most
$c$, one must also impose $\chi\le c$.  The condition
$\chi\le t_r$ alone does not preserve such a third-slot bound.

In the classical case $\nu<1/2$ with a fixed margin, use the absolute
Bombieri--Vinogradov remainder for complete moduli.  There is no
full-well or ordering requirement on the specified prefix in that
case.  Each complete modulus has only boundedly many ways to contain
the finitely many labelled prime slots, because every prime slot
exceeds a fixed positive power of $N$.  The residual construction
has bounded coefficients for each fixed precision, so the complete
coefficient remains bounded.  This is the classical interface of
Section~\ref{sec:generalized-one-step-double-sieve}, including its
coprimality and collision treatment.

\begin{proposition}[Second seed with stable prefix data]
\label{prop:second-seed-stable-prefix}
Under the named moving-residue and positive $K=681/100$ inputs of
the preceding seed, suppose $\alpha\ge2$ and one of these interfaces
holds.  In the enhanced case assume additionally
\begin{equation}\label{eq:second-seed-comparison-cutoff}
 L/s_*<\chi.
\end{equation}
Then, with the same normalization $\Theta_Q$, one has
\[
 \Phi_{s_*}\le(1-J_*+o(1))\Theta_Q.
\]
For any $0<w\le s_*$, positivity and cutoff monotonicity imply the
same upper bound for $\Phi_w$.  In particular the normalized ordinary
upper multiplier $A(w)=1$ for $1\le w\le s_*$ may be replaced by
$1-J_*$.
\end{proposition}

\begin{proof}
All specified primes added by the 21-term inequality are below
$R^{1/s_*}$.  For every non-switched term the subsequently sieved
prime slots are smaller still.  Under
\eqref{eq:second-seed-comparison-cutoff} all these slots are below
the old specified prefix and below its stable ceiling $\chi$.
Thus the original largest prime, original selected subset, its old
remaining-square gates, and the applicable $W$ data remain unchanged.

The residual singleton in each negative term satisfies $2u<1$,
where $u\le1/s_*<1/2$ is relative to $R$.  The residual pair in
$\Gamma_6,\Gamma_7,\Gamma_8$ satisfies both $2v<1$ and $v+2u<1$.
These same tests hold on the good portion of $\Gamma_5$.
The exact merge lemma proved with the seed combines each such
residual vector with the old full-well prefix.  It then combines
their ordinary residual sieve atoms at the correct residual level.
Equivalently, since all old slots precede all new slots, the old
complete and selected-remaining tests are unchanged, and each new
test is the corresponding residual square test multiplied by $B$.
This explicitly checks the complete atom, not just its bare prefix.

Every factorization is performed on buffered endpoint vectors.
Keep all marked slots, the original scalar signs, and the original
residual factorials; use the existing bounded-group normalization
and moving-residue argument.  The small-prime weight and the
exclusion of primes dividing the specified prefix are treated by
the same signed interface.  Thus all non-switched ordinary terms
retain exactly their previous main kernels.

On the bad portion of $\Gamma_5$ no signed distribution statement
is used.  The positive product estimate gives the existing
replacement $3Ks_0/20$.  The proof of every positive switched mass
uses only $\alpha\ge2$, a fixed positive exponent cutoff, and the
retained labelled domains.  It is independent of the former anchor
interval and of the largest exponent of $B$.  In particular the
uniform $M_{20},M_{21}$ still include their separate unit surfaces;
one does not reuse their special $\alpha=2$ simplifications.

In the classical case, all complete ordinary moduli have support
at most $Q$, including on the bad portion of $\Gamma_5$.  Absolute
Bombieri--Vinogradov therefore licenses even that ordinary kernel.
For a common conservative formula we may instead keep the positive
replacement there; either choice is a valid upper bound.  No
selected-subproduct or full-well assertion is needed for this case.

All main kernels in the certified functional are consequently
unchanged.  Its already saved outward rational certificate gives
$J_2>9/500$.  Finally, if $w\le s_*$, then
$R^{1/w}\ge R^{1/s_*}$ and
$S(\mathcal A_{N,B},R^{1/w})\le
S(\mathcal A_{N,B},R^{1/s_*})$ pointwise.  The comparison cutoff
is lowered; the normalization uses the same total level $Q$.
\end{proof}

\paragraph{The correction to use and the exact lower transfer.}
Write
\[
 j_2(w)=J_*\mathbf1_{1\le w\le s_*}.
\]
At an upper node with $1\le w\le s_*$ and physical cutoff
$R^{1/w}\le N^\chi$, condition
\eqref{eq:second-seed-comparison-cutoff} follows automatically,
after the usual trimming of equality sets.  More generally it is
enough to check the comparison cutoff directly.  If $j_1$ is any
proved first-seed correction on the same physical prefix box, use
$\max\{j_1,j_2\}$ there.  Their defects are alternative upper
bounds for the same count and are not added.

For a stable lower prefix with cutoff $N^a\le N^\chi$, write
$s=L/a\ge2$.  Buchstab from $R^{1/4}$ to $R^{1/s}$, when $s<4$,
gives children indexed by $w\in[s-1,3]$.  The appended prime and
the child ratio are
\[
 q=N^{L/(w+1)},\quad L'=\frac{Lw}{w+1},\qquad
 \alpha'=\alpha+\frac{\alpha-1}{w}\ge2.
\]
Every child used by the second seed has $1\le w\le s_*$.
Its new square condition follows from $w\ge1$, namely
$2\log_Nq\le L$.  Also $q\le N^a\le N^\chi$, so the original
branch persists, and its seed comparison cutoff satisfies
\begin{equation}\label{eq:second-seed-child-cutoff}
 \frac{L'}{s_*}=\frac{L}{w+1}\frac{w}{s_*}
 \le\log_Nq.
\end{equation}
In particular the seed's new prime slots are below this child's
new smallest specified prime.  This remains true for $1\le w<2$.
It would lose a valid part of the estimate to truncate the integral
at $w=2$.

The ordinary identity
$a(4)-\int_{s-1}^3 A(w)\,dw/w=a(s)$ holds for $2\le s<4$,
including $A(w)=1$ on $1\le w<2$.  After inserting the uniform seed
defect and retaining its exact prefix measure $dw/w$, the new lower
correction is
\begin{equation}\label{eq:second-seed-stable-lower}
 h_2(s)=
 \begin{cases}
 \displaystyle\frac9{500}\log\frac{11/5}{s-1},&2\le s<16/5,\\
 0,&s\ge16/5.
 \end{cases}
\end{equation}
Thus the normalized lower multiplier is at least $a(s)+h_2(s)$.
This is a lower bound from one finite tree; it is not an abstract
old $h$ value.  At $s=2$, its limiting value is
$(9/500)\log(11/5)$, approximately $0.01419$.
When combining seeds, integrate their pointwise maximum on the
child boxes; do not add $h_2$ to a first-seed correction that already
uses any of the same children.

\paragraph{Actual domains in the complete eleven-term count.}
Use the parameters and level prescriptions of
Section~\ref{sec:generalized-one-step-double-sieve}.  All the
following assertions concern the limiting cap approached from
strictly admissible levels.
\begin{itemize}
\item For $L_2$, the largest-prime branch has
$\nu=(1+t)/2$ and target $w=(1-t)/(2t)$, on the relevant
$t<h<25/128$ range.  The direct second seed applies exactly when
$5/27\le t<h$.  This interval is empty for the current sampled
$h\simeq0.1505$; the separately proved finite $H(3)$ propagation
on $0.145<t<0.15$ is a different available option.
\item For $U_3,U_4$, set $L=\nu-t$ and use the current physical
upper argument $w=\max\{2,L/a\}$.  The direct seed is applicable
on $0<L\le(11/5)a$, since its comparison cutoff is at most $a<t$.
Here $\nu=v(t)$ on the enhanced part.  On the classical $U_3$
tail $57/224<t<1/3$, it applies on
$t\ge1/2-(11/5)a$.  Keep the corrected coefficient
$4(1-\max\{j_1,j_2\})/L$ in place of the ordinary option.
\item For $L_5,L_6$, retain their existing level $V_{2,a}(t,u)$
and all its gates.  Use \eqref{eq:second-seed-stable-lower} where
$2\le(\nu-t-u)/a<16/5$.  In this application the first appended
prime $q$ is at most $N^a$, so the two original primes remain the
largest.  If the chosen $W$ cap used a third-slot bound $a$, this
child now has an actual third slot $q\le N^a$.  Every seed prime
is below $q$ by \eqref{eq:second-seed-child-cutoff}, and the bound
is therefore preserved.  The selected leading pair is preserved
in the same way.  The gain integrand is
$h_2((\nu-t-u)/a)/(tu(\nu-t-u))$ before combining with any
overlapping first-seed correction.
\item For $U_{10},U_{11}$, retain the full eligibility set and
all branch gates of the four-prefix safe minimum.  Put
$T=x+y+z+w_0$, $L=\nu-T$ and let $w_0$ be the smallest specified
exponent.  The physical upper argument is
$\max\{2,L/w_0\}$.  The direct second seed applies when
$L\le(11/5)w_0$, and its comparison cutoff is then below the
smallest old slot.  All original selected slots and the first
three $W$ slots persist.  Replace the ordinary option inside the
existing switching/ordinary minimum by
$4(1-\max\{j_1,j_2\})/L$; do not add this saving to an already
overlapping switched improvement.
\end{itemize}
For a common finite positive partition these are genuine averaged
prefix bounds.  The endpoint comparisons and cutoff equality sets
are approached through buffered boxes.  There is no assertion for
an individual fixed tuple of primes.

\paragraph{Finite recursive closure, not a restored old table.}
The state conditions also permit further finite Buchstab trees.
At an upper or lower parent with physical argument $s\ge2$ and
cutoff below its stable ceiling, an appended prime indexed by
$w\ge s-1\ge1$ has square at most the current residual level and
is below that ceiling.  Retaining the original branch data and
using this new prime as the child's cutoff preserves the same
typed family.  Its actual ratio changes to
$\alpha+(\alpha-1)/w$; it is not reset to two.

For example, if uniform defect bounds $H,h$ have already been proved
on all nodes of a finite such tree, and $S_{\max}+1\ge s$, ordinary
Buchstab cancellation gives the further valid choices
\[
 h_{\rm new}(s)=\int_{s-1}^{S_{\max}}H(w)\frac{dw}{w},\qquad
 H_{\rm new}(s)=\int_{s-1}^{S_{\max}}h(w)\frac{dw}{w},
 \quad s\ge2,
\]
where unproved corrections are set to zero; in particular $h(w)=0$
for $w<2$, whereas the proved seed gives $H(w)=J_*$ for
$1\le w\le s_*$.  These are alternative lower bounds for the
defect, so one may take their maximum with a previous bound rather
than adding overlapping estimates.  Every generation is finite,
and its positive prime cutoff remains fixed before $N\to\infty$.
The first- and second-seed identities can similarly use previously
proved child defects at their actual states, with their original
coefficients.  This paragraph licenses an explicitly checked finite
construction; it does not supply an unevaluated numerical matrix
or an infinite-iteration convergence claim.

One explicit consequence needs no new integration certificate.
If the upper comparison cutoff $R^{1/3}$ is below the stable ceiling,
the same two Buchstab steps used for
\eqref{eq:replacement-second-seed-h3}, now in this enlarged family,
give
\begin{equation}\label{eq:second-seed-stable-h3}
 H(3)\ge J_*\log(5/4)\log(22/15)>7/4625.
\end{equation}
The first new prime is at most $R^{1/3}$; the second is at most
the first, and the final seed comparison cutoff is at most the
second.  Thus every old branch datum persists.  The proof uses
only the uniform $\alpha\ge2$ seed, so this consequence has no
restriction to the original anchor interval.  By count monotonicity
it also bounds upper nodes with physical argument $1\le w\le3$,
provided the comparison cutoff $R^{1/3}$ is licensed.  It is another
option to maximize with the direct seed and other proved defects.

The saved companion JSON binds the already certified scalar seed
and records these analytic domains.  No new total counting
coefficient, transfer exponent, global optimality claim, or
resolution of the requested frontier follows from this extension
alone.  Any numerical improvement after taking the common-box
maxima requires its own outward integral certificate.
